\section{Discussion}
\label{sec:discussion}

% Intended content (~2 pages).
%
%  8.1 What the geometry buys.
%      Summarise: a coordinate-free description of paired efficiency in which
%      the four-quadrant taxonomy of the empirical companion is a coordinate
%      artefact; a distribution-free signal/information pair via the
%      partition-function reading; a meta-analytic scale (psi) on which
%      cross-domain effect sizes are properly comparable; a single
%      Riemann-sphere figure that organises the framework.
%
%  8.2 Information geometry context.
%      The paired-efficiency family sits inside the broader information-
%      geometric program initiated by Rao (1945), Amari (1985), and others.
%      Brief comparison with Fisher information metrics on more familiar
%      exponential families.  This is the section that earns the paper
%      its place in the information-geometry literature.
%
%  8.3 Relation to the empirical companion paper.
%      Concretely: which results in this paper are *used* by the empirical
%      companion (the kappa <-> 1/kappa symmetry, the psi coordinate, the
%      partition-function remark, the regime-change observation), and which
%      are stand-alone contributions of this paper (the canonical form,
%      Moebius linearisation, Chebyshev/Poisson expansion, sphere lift,
%      Fisher--Rao identification).
%
%  8.4 Limitations.
%      - All results are univariate-pair: an extension to multivariate paired
%        efficiency (PEF matrices) is outlined in Section 8.5 but not
%        developed here.
%      - The partition-function reading requires rho sech tau in (0, 1);
%        the rho < 0 regime is mathematically distinct and is treated by
%        sign-carry, not by analytic continuation.
%      - Numerical validation depends on bootstrap inference; finite-sample
%        coverage of the psi-scale intervals is checked but not theoretically
%        derived to high order.
%
%  8.5 Future directions.
%      Six concrete extensions:
%        (i)   multivariate PEF: a matrix-valued generalisation with its own
%              spectral / spherical structure;
%        (ii)  non-parametric psi: substitute rank correlation for rho-hat
%              and derive the corresponding stabilising transform;
%        (iii) operator-theoretic / Hardy-space treatment building on the
%              Poisson-kernel decomposition of Section 3.4;
%        (iv)  small-sample bias correction for eta-hat and psi-hat;
%        (v)   regime-change asymptotics: refined treatment of the rho = 0
%              singularity in the latent geometric family;
%        (vi)  causal-inference connections: when does correlation-based
%              variance reduction preserve causal estimates?
%
%  8.6 Closing.
%      One-paragraph summary: Fisher's generalised paired efficiency formula
%      is, beneath the algebra, a piece of spherical and information geometry
%      with a distribution-free signal/information pair.  The empirical
%      companion paper validates the formula across six domains; this paper
%      explains *why* the validation works as well as it does, and provides
%      the meta-analytic infrastructure for future cross-domain pooling.

% TODO: drafting.
